% GENERATED FILE. Edit canonical lessons, then run npm run generate:books.
% canonical-source-hash: fnv1a64:f4fb579119b3e7be
% canonical-lessons: FR-C316-celebrer, FR-C316-celebration, FR-C316-cortege, FR-C316-fanfare, FR-C316-jour-ferie

\chapter{Célébrer}
\label{ch:fr-celebrer}
\hlchaptermodality{316}

\begin{chapteropening}
\textbf{By the end of this chapter:} \emph{I can describe a public celebration: celebrating, the celebration, a procession, a brass band and a public holiday.}

Complete the last lesson of chapter 316: I can describe a public celebration: celebrating, the celebration, a procession, a brass band and a public holiday.
\end{chapteropening}

\section[célébrer]{célébrer — to celebrate}

\label{lesson:FR-C316-celebrer}



\begin{quote}
Before the new one: say the French for to inherit, then the French for the inheritance.
\end{quote}

\subsection*{You'll want to know: célébrer}
\textbf{célébrer} — \textquotedblleft{}to celebrate\textquotedblright{}.

\begin{quote}
La ville a célébré la fin de la guerre.
\end{quote}

\begin{quote}
\emph{The town celebrated the end of the war.}
\end{quote}

More formal than \emph{fêter}; also for a mass or a wedding: \emph{célébrer un mariage}.

\begin{cousinweb}
\emph{célébrer} is a learned borrowing from Latin \emph{celebrāre} \textquotedblleft{}to frequent, to celebrate, to honour\textquotedblright{} (EtymDB), built on \emph{celeber} \textquotedblleft{}crowded, famous\textquotedblright{}.
\end{cousinweb}

\begin{grammarlens}[title={What you've built}]
One more word for telling, explaining and describing at B1.
\end{grammarlens}

\subsection*{Your turn}
\begin{itemize}
\raggedright
  \item \emph{Say it:} \emph{célébrer}
  \item \emph{Say it:} \emph{célébrer}, once more
  \item \emph{Say it:} a sentence of your own with \emph{célébrer}
  \item \emph{Recall:} say the French for to inherit, then the French for the inheritance, then say \emph{célébrer} again
\end{itemize}

\begin{tcolorbox}[breakable,colback=teal!4,colframe=teal!35!black,title={Before you move on}]
What does célébrer mean? (\textquotedblleft{}To celebrate\textquotedblright{}.) Say it once more.
\end{tcolorbox}
\section[la célébration]{la célébration — the celebration}

\label{lesson:FR-C316-celebration}



\begin{quote}
Before the new one: say the French for the inheritance, then the French for to celebrate.
\end{quote}

\subsection*{You'll want to know: la célébration}
\textbf{la célébration} — \textquotedblleft{}the celebration\textquotedblright{}.

\begin{quote}
La célébration a duré toute la nuit.
\end{quote}

\begin{quote}
\emph{The celebration went on all night.}
\end{quote}

Often for a religious or official event.

\begin{cousinweb}
\emph{célébration} is a learned borrowing from Latin \emph{celebrātiō} \textquotedblleft{}a crowd, a festival\textquotedblright{}, from \emph{celebrāre} (Lewis \& Short).
\end{cousinweb}

\begin{grammarlens}[title={What you've built}]
One more word for telling, explaining and describing at B1.
\end{grammarlens}

\subsection*{Your turn}
\begin{itemize}
\raggedright
  \item \emph{Say it:} \emph{la célébration}
  \item \emph{Say it:} \emph{la célébration}, once more
  \item \emph{Say it:} a sentence of your own with \emph{la célébration}
  \item \emph{Recall:} say the French for the inheritance, then the French for to celebrate, then say \emph{la célébration} again
\end{itemize}

\begin{tcolorbox}[breakable,colback=teal!4,colframe=teal!35!black,title={Before you move on}]
What does la célébration mean? (\textquotedblleft{}The celebration\textquotedblright{}.) Say it once more.
\end{tcolorbox}
\section[le cortège]{le cortège — the procession}

\label{lesson:FR-C316-cortege}



\begin{quote}
Before the new one: say the French for to celebrate, then the French for the celebration.
\end{quote}

\subsection*{You'll want to know: le cortège}
\textbf{le cortège} — \textquotedblleft{}the procession\textquotedblright{}.

\begin{quote}
Le cortège a traversé le centre-ville.
\end{quote}

\begin{quote}
\emph{The procession went through the town centre.}
\end{quote}

A wedding, funeral or protest procession.

\begin{cousinweb}
\emph{cortège} was borrowed from Italian \emph{corteggio} \textquotedblleft{}a retinue\textquotedblright{} (EtymDB), built on \emph{corte} \textquotedblleft{}court\textquotedblright{}, from Latin \emph{cohors} \textquotedblleft{}an enclosure, a company\textquotedblright{} (Lewis \& Short).
\end{cousinweb}

\begin{grammarlens}[title={What you've built}]
One more word for telling, explaining and describing at B1.
\end{grammarlens}

\subsection*{Your turn}
\begin{itemize}
\raggedright
  \item \emph{Say it:} \emph{le cortège}
  \item \emph{Say it:} \emph{le cortège}, once more
  \item \emph{Say it:} a sentence of your own with \emph{le cortège}
  \item \emph{Recall:} say the French for to celebrate, then the French for the celebration, then say \emph{le cortège} again
\end{itemize}

\begin{tcolorbox}[breakable,colback=teal!4,colframe=teal!35!black,title={Before you move on}]
What does le cortège mean? (\textquotedblleft{}The procession\textquotedblright{}.) Say it once more.
\end{tcolorbox}
\section[la fanfare]{la fanfare — the brass band; a fanfare}

\label{lesson:FR-C316-fanfare}



\begin{quote}
Before the new one: say the French for the celebration, then the French for the procession.
\end{quote}

\subsection*{You'll want to know: la fanfare}
\textbf{la fanfare} — \textquotedblleft{}the brass band; a fanfare\textquotedblright{}.

\begin{quote}
La fanfare jouait devant la mairie.
\end{quote}

\begin{quote}
\emph{The brass band was playing in front of the town hall.}
\end{quote}

\emph{En fanfare} — with a lot of noise.

\begin{cousinweb}
\emph{fanfare} is usually explained as imitating the sound of trumpets; the origin is not settled (EtymDB lists no parent).
\end{cousinweb}

\begin{grammarlens}[title={What you've built}]
One more word for telling, explaining and describing at B1.
\end{grammarlens}

\subsection*{Your turn}
\begin{itemize}
\raggedright
  \item \emph{Say it:} \emph{la fanfare}
  \item \emph{Say it:} \emph{la fanfare}, once more
  \item \emph{Say it:} a sentence of your own with \emph{la fanfare}
  \item \emph{Recall:} say the French for the celebration, then the French for the procession, then say \emph{la fanfare} again
\end{itemize}

\begin{tcolorbox}[breakable,colback=teal!4,colframe=teal!35!black,title={Before you move on}]
What does la fanfare mean? (\textquotedblleft{}The brass band; a fanfare\textquotedblright{}.) Say it once more.
\end{tcolorbox}
\section[le jour férié]{le jour férié — the public holiday, the bank holiday}

\label{lesson:FR-C316-jour-ferie}



\begin{quote}
Before the new one: say the French for the procession, then the French for the brass band.
\end{quote}

\subsection*{You'll want to know: le jour férié}
\textbf{le jour férié} — \textquotedblleft{}the public holiday, the bank holiday\textquotedblright{}.

\begin{quote}
Le premier mai est un jour férié.
\end{quote}

\begin{quote}
\emph{The first of May is a public holiday.}
\end{quote}

\emph{Faire le pont} — to take the day off between a holiday and the weekend.

\begin{cousinweb}
\emph{férié} is a learned borrowing from Latin \emph{fēriātus} \textquotedblleft{}on holiday\textquotedblright{}, from \emph{fēriae} \textquotedblleft{}days of rest, holidays\textquotedblright{} (Lewis \& Short); \emph{jour} comes from Latin \emph{diurnum}.
\end{cousinweb}

\begin{grammarlens}[title={What you've built}]
One more word for telling, explaining and describing at B1.
\end{grammarlens}

\subsection*{Your turn}
\begin{itemize}
\raggedright
  \item \emph{Say it:} \emph{le jour férié}
  \item \emph{Say it:} \emph{le jour férié}, once more
  \item \emph{Say it:} a sentence of your own with \emph{le jour férié}
  \item \emph{Recall:} say the French for the procession, then the French for the brass band, then say \emph{le jour férié} again
\end{itemize}

\begin{tcolorbox}[breakable,colback=teal!4,colframe=teal!35!black,title={Before you move on}]
What does le jour férié mean? (\textquotedblleft{}The public holiday, the bank holiday\textquotedblright{}.) Say it once more.
\end{tcolorbox}
